%%%%%%%%%%%%%%%%%%%%%%%%%%%%%%%%%%%%%%%%%%%%%%%%%%%%%%%%%%%%%%%%%
% PHIẾU HỌC TẬP - EPIDEMICSIM: MÔ PHỎNG LÂY LAN DỊCH BỆNH
%%%%%%%%%%%%%%%%%%%%%%%%%%%%%%%%%%%%%%%%%%%%%%%%%%%%%%%%%%%%%%%%%
\documentclass[12pt,a4paper]{article}

\usepackage[utf8]{inputenc}
\usepackage[T1]{fontenc}
\usepackage[vietnamese]{babel}
\usepackage{amsmath,amssymb}
\usepackage{array}
\usepackage{booktabs}
\usepackage{tabularx}
\usepackage{colortbl}
\usepackage{graphicx}
\usepackage{newtxtext, newtxmath}
\usepackage[scaled=0.92]{helvet}
\usepackage{microtype}

\usepackage[svgnames]{xcolor}
\definecolor{primary}{RGB}{26,95,122}
\definecolor{accent}{RGB}{255,107,53}
\definecolor{danger}{RGB}{194,55,40}
\definecolor{success}{RGB}{87,204,153}
\definecolor{TableHeader}{gray}{0.9}

\usepackage[left=1.5cm,right=1.5cm,top=2cm,bottom=2cm,headheight=15pt]{geometry}
\usepackage{fancyhdr}
\pagestyle{fancy}
\fancyhf{}
\renewcommand{\headrule}{\color{primary}\hrule width\textwidth height 1pt}
\fancyhead[L]{\color{primary}\footnotesize\sffamily\textbf{PHIẾU HỌC TẬP - EPIDEMICSIM}}
\fancyhead[R]{\color{primary}\footnotesize\sffamily\textbf{LỚP 11}}
\fancyfoot[C]{\sffamily\thepage}

\usepackage{enumitem}
\usepackage{tcolorbox}
\tcbuselibrary{skins, breakable}
\usepackage[colorlinks=true,linkcolor=primary,urlcolor=accent]{hyperref}

\newtcolorbox{phieuhoctap}[1]{
    enhanced, breakable,
    colback=white, colframe=danger,
    boxrule=1.5pt, arc=3mm,
    fonttitle=\bfseries\sffamily\Large,
    coltitle=white, colbacktitle=danger,
    title=#1
}

\newtcolorbox{tinhuong}{
    enhanced, breakable,
    colback=primary!5, colframe=primary!80,
    boxrule=1pt, arc=2mm,
    fonttitle=\bfseries\sffamily,
    title=\color{accent}TÌNH HUỐNG THỰC TIỄN
}

\newtcolorbox{congthuc}{
    enhanced, breakable,
    colback=accent!5, colframe=accent!80,
    boxrule=1pt, arc=2mm,
    fonttitle=\bfseries\sffamily,
    title=\color{primary}CÔNG THỨC TRỌNG TÂM
}

\newtcolorbox{khaosat}{
    enhanced, breakable,
    colback=gray!5, colframe=gray!60,
    boxrule=1pt, arc=2mm,
    fonttitle=\bfseries\sffamily,
    title=\color{primary}PHIẾU KHẢO SÁT
}

\newcolumntype{C}{>{\centering\arraybackslash}X}
\newcolumntype{L}{>{\raggedright\arraybackslash}X}
\newcommand{\header}[1]{\cellcolor{TableHeader}\sffamily\bfseries#1}

\setlist[itemize,1]{label=\textbullet, leftmargin=*, nosep}
\setlist[itemize,2]{label=--, leftmargin=*, nosep}

\begin{document}

%==============================================================
% PHIẾU HỌC TẬP SỐ 1
%==============================================================
\begin{phieuhoctap}{PHIẾU HỌC TẬP SỐ 1: KHÁM PHÁ MÔ HÌNH SIR}

\noindent\textbf{Họ và tên:} \underline{\hspace{6cm}} \hfill \textbf{Nhóm:} \underline{\hspace{2cm}}

\vspace{0.5cm}
\begin{tinhuong}
Năm 2020, đại dịch COVID-19 bùng phát. Tại một thành phố có \textbf{10,000 dân}, ban đầu có \textbf{10 ca nhiễm}. Các nhà khoa học sử dụng mô hình SIR để dự đoán diễn biến dịch.

\textbf{Thông số:}
\begin{itemize}
\item Tỷ lệ lây nhiễm: $\beta = 0.3$ (mỗi người nhiễm lây cho 0.3 người/ngày)
\item Tỷ lệ hồi phục: $\gamma = 0.1$ (10\% người nhiễm khỏi bệnh mỗi ngày)
\end{itemize}
\end{tinhuong}

\vspace{0.5cm}
\textbf{\color{accent}BÀI 1: Xác định 3 nhóm dân số SIR}
\begin{itemize}
\item \textbf{S} (Susceptible - Nhạy cảm): Số người có thể bị lây = \underline{\hspace{4cm}} người
\item \textbf{I} (Infected - Đang nhiễm): Số người đang nhiễm bệnh = \underline{\hspace{4cm}} người
\item \textbf{R} (Recovered - Hồi phục): Số người đã khỏi bệnh = \underline{\hspace{4cm}} người
\item \textbf{Kiểm tra:} S + I + R = \underline{\hspace{2cm}} = N (dân số)
\end{itemize}

\vspace{0.5cm}
\textbf{\color{accent}BÀI 2: Tính hệ số lây nhiễm cơ bản $R_0$}
\begin{itemize}
\item Công thức: $R_0 = \dfrac{\beta}{\gamma}$
\item Thay số: $R_0 = \dfrac{\underline{\hspace{1.5cm}}}{\underline{\hspace{1.5cm}}} = \underline{\hspace{2cm}}$
\item \textbf{Ý nghĩa:} Mỗi người nhiễm trung bình lây cho \underline{\hspace{2cm}} người khác
\item Nếu $R_0 > 1$: \underline{\hspace{6cm}} (dịch bùng phát / dịch tự tắt)
\end{itemize}

\vspace{0.5cm}
\textbf{\color{accent}BÀI 3: Tính ngưỡng miễn dịch cộng đồng}
\begin{itemize}
\item Công thức: $H = 1 - \dfrac{1}{R_0}$
\item Thay số: $H = 1 - \dfrac{1}{\underline{\hspace{1.5cm}}} = \underline{\hspace{2cm}}$ = \underline{\hspace{2cm}}\%
\item \textbf{Ý nghĩa:} Cần \underline{\hspace{2cm}}\% dân số có miễn dịch để dịch tự tắt
\end{itemize}

\vspace{0.5cm}
\textbf{\color{accent}BÀI 4: Hàm số mũ trong giai đoạn đầu}
\begin{itemize}
\item Số ca nhiễm theo thời gian: $I(t) = I_0 \cdot e^{(\beta - \gamma) \cdot t}$
\item Với $I_0 = 10$, $\beta = 0.3$, $\gamma = 0.1$:
\item $I(10) = 10 \cdot e^{(0.3 - 0.1) \times 10} = 10 \cdot e^{\underline{\hspace{1cm}}} = \underline{\hspace{4cm}}$ ca
\item $I(20) = 10 \cdot e^{(0.3 - 0.1) \times 20} = 10 \cdot e^{\underline{\hspace{1cm}}} = \underline{\hspace{4cm}}$ ca
\end{itemize}

\begin{congthuc}
\textbf{Mô hình SIR:}
$$\frac{dS}{dt} = -\beta \cdot \frac{S \cdot I}{N}, \quad \frac{dI}{dt} = \beta \cdot \frac{S \cdot I}{N} - \gamma \cdot I, \quad \frac{dR}{dt} = \gamma \cdot I$$

\textbf{Hệ số lây nhiễm:} $R_0 = \dfrac{\beta}{\gamma}$ \hspace{2cm} \textbf{Ngưỡng miễn dịch:} $H = 1 - \dfrac{1}{R_0}$
\end{congthuc}
\end{phieuhoctap}

\newpage
%==============================================================
% PHIẾU HỌC TẬP SỐ 2
%==============================================================
\begin{phieuhoctap}{PHIẾU HỌC TẬP SỐ 2: SỬ DỤNG EPIDEMICSIM}

\noindent\textbf{Họ và tên:} \underline{\hspace{6cm}} \hfill \textbf{Nhóm:} \underline{\hspace{2cm}}

\vspace{0.5cm}
\textbf{\color{primary}I. GIỚI THIỆU ỨNG DỤNG}
\begin{itemize}
\item \textbf{Tên ứng dụng:} EpidemicSim - Mô phỏng lây lan dịch bệnh
\item \textbf{Địa chỉ:} \href{https://stem-1h8.pages.dev/EpidemicSim.html}{stem-1h8.pages.dev/EpidemicSim.html} hoặc file \texttt{EpidemicSim.html}
\item \textbf{Cấu trúc 3 Tab:}
    \begin{itemize}
    \item \textbf{Tab 1 - MÔ PHỎNG:} Nhập tham số, xem biểu đồ SIR động
    \item \textbf{Tab 2 - SO SÁNH KỊCH BẢN:} Thêm kịch bản, so sánh đỉnh dịch
    \item \textbf{Tab 3 - HỌC TẬP:} Công thức, kiến thức dịch tễ học
    \end{itemize}
\end{itemize}

\vspace{0.5cm}
\textbf{\color{primary}II. NHIỆM VỤ THỰC HÀNH}

\textbf{\color{accent}Bước 1: Tab ``MÔ PHỎNG'' - Kịch bản gốc}
\begin{itemize}
\item Nhập: N = 10,000 | $I_0$ = 10 | $\beta$ = 0.3 | $\gamma$ = 0.1 | Tiêm chủng = 0\%
\item Nhấn ``Chạy mô phỏng''
\end{itemize}

\textbf{Ghi lại kết quả:}

\begin{tabularx}{\textwidth}{|L|C|C|C|}
\hline
& \header{Đỉnh dịch (số ca max)} & \header{Ngày đỉnh dịch} & \header{Tổng số ca nhiễm} \\
\hline
Kịch bản 1 (không can thiệp) & \underline{\hspace{2cm}} & \underline{\hspace{2cm}} & \underline{\hspace{2cm}} \\
\hline
\end{tabularx}

\vspace{0.5cm}
\textbf{\color{accent}Bước 2: Tab ``SO SÁNH KỊCH BẢN''}
\begin{itemize}
\item Thêm \textbf{Kịch bản 2 - Giãn cách xã hội}: $\beta$ = 0.2 (giảm tiếp xúc)
\item Thêm \textbf{Kịch bản 3 - Tiêm chủng 50\%}: Giữ nguyên $\beta$, Tiêm chủng = 50\%
\item Nhấn ``So sánh''
\end{itemize}

\textbf{Ghi lại kết quả so sánh:}

\begin{tabularx}{\textwidth}{|L|C|C|}
\hline
\header{Kịch bản} & \header{Đỉnh dịch (số ca)} & \header{Hiệu quả so với KB1} \\
\hline
Kịch bản 1 - Không can thiệp & \underline{\hspace{3cm}} & -- \\
\hline
Kịch bản 2 - Giãn cách xã hội & \underline{\hspace{3cm}} & Giảm \underline{\hspace{2cm}}\% \\
\hline
Kịch bản 3 - Tiêm chủng 50\% & \underline{\hspace{3cm}} & Giảm \underline{\hspace{2cm}}\% \\
\hline
\end{tabularx}

\vspace{0.5cm}
\textbf{\color{accent}Bước 3: Tab ``HỌC TẬP''}
\begin{itemize}
\item Đọc phần ``Ngưỡng miễn dịch cộng đồng''
\item Với COVID-19 ($R_0 \approx 2.5$): Cần tiêm chủng \underline{\hspace{2cm}}\% dân số
\item Với Sởi ($R_0 \approx 15$): Cần tiêm chủng \underline{\hspace{2cm}}\% dân số
\end{itemize}

\vspace{0.5cm}
\textbf{\color{primary}III. KẾT LUẬN}

\textbf{Kịch bản tốt nhất:} \underline{\hspace{8cm}}

\textbf{Chính sách phòng dịch nhóm em đề xuất:}

\underline{\hspace{\textwidth}}

\underline{\hspace{\textwidth}}
\end{phieuhoctap}

\newpage
%==============================================================
% PHIẾU KHẢO SÁT TRƯỚC KHI HỌC
%==============================================================
\begin{khaosat}

\begin{center}
\Large\bfseries\sffamily PHIẾU KHẢO SÁT TRƯỚC KHI HỌC
\end{center}

\noindent\textbf{Họ và tên:} \underline{\hspace{6cm}} \hfill \textbf{Lớp:} \underline{\hspace{2cm}}

\vspace{0.5cm}
\textbf{Đánh dấu ($\checkmark$) vào ô phù hợp:}

\vspace{0.3cm}
\begin{tabularx}{\textwidth}{|L|c|c|c|c|}
\hline
\header{Câu hỏi} & \header{5} & \header{4} & \header{3} & \header{1-2} \\
\hline
1. Em hiểu dịch bệnh lây lan như thế nào & $\square$ & $\square$ & $\square$ & $\square$ \\
\hline
2. Em biết $R_0$ (hệ số lây nhiễm) là gì & $\square$ & $\square$ & $\square$ & $\square$ \\
\hline
3. Em hiểu ``miễn dịch cộng đồng'' là gì & $\square$ & $\square$ & $\square$ & $\square$ \\
\hline
4. Em biết Hàm số mũ ứng dụng trong y tế & $\square$ & $\square$ & $\square$ & $\square$ \\
\hline
5. Em thấy Toán giúp dự đoán diễn biến dịch & $\square$ & $\square$ & $\square$ & $\square$ \\
\hline
\end{tabularx}

\textit{\footnotesize Thang điểm: 5=Rất đồng ý, 4=Đồng ý, 3=Phân vân, 1-2=Không đồng ý}

\vspace{0.5cm}
\textbf{6. Em mong muốn học được gì từ bài học này?}

\underline{\hspace{\textwidth}}
\end{khaosat}

\vspace{1cm}
%==============================================================
% PHIẾU KHẢO SÁT SAU KHI HỌC
%==============================================================
\begin{khaosat}

\begin{center}
\Large\bfseries\sffamily PHIẾU KHẢO SÁT SAU KHI HỌC
\end{center}

\noindent\textbf{Họ và tên:} \underline{\hspace{6cm}} \hfill \textbf{Lớp:} \underline{\hspace{2cm}}

\vspace{0.5cm}
\begin{tabularx}{\textwidth}{|L|c|c|c|c|}
\hline
\header{Câu hỏi} & \header{5} & \header{4} & \header{3} & \header{1-2} \\
\hline
1. Em hiểu mô hình SIR và ý nghĩa của $R_0$ & $\square$ & $\square$ & $\square$ & $\square$ \\
\hline
2. Em biết cách tính ngưỡng miễn dịch cộng đồng & $\square$ & $\square$ & $\square$ & $\square$ \\
\hline
3. Ứng dụng EpidemicSim giúp em hiểu bài tốt hơn & $\square$ & $\square$ & $\square$ & $\square$ \\
\hline
4. Em thấy Toán học hữu ích trong phòng chống dịch & $\square$ & $\square$ & $\square$ & $\square$ \\
\hline
5. Em tự tin giải thích được tại sao cần tiêm chủng & $\square$ & $\square$ & $\square$ & $\square$ \\
\hline
\end{tabularx}

\textit{\footnotesize Thang điểm: 5=Rất đồng ý, 4=Đồng ý, 3=Phân vân, 1-2=Không đồng ý}

\vspace{0.5cm}
\textbf{6. Điều em thích nhất về bài học này là gì?}

\underline{\hspace{\textwidth}}

\textbf{7. Em muốn cải thiện gì về bài học này?}

\underline{\hspace{\textwidth}}
\end{khaosat}

\newpage
%==============================================================
% RUBRIC ĐÁNH GIÁ
%==============================================================
\begin{center}
{\Large\bfseries\sffamily\color{primary} RUBRIC ĐÁNH GIÁ SẢN PHẨM}
\end{center}

\vspace{0.3cm}
\noindent\textbf{Tên nhóm:} \underline{\hspace{4cm}} \hfill \textbf{Người đánh giá:} \underline{\hspace{4cm}}

\vspace{0.5cm}
\begin{tabularx}{\textwidth}{|>{\raggedright\arraybackslash}p{2.5cm}|>{\raggedright\arraybackslash}X|>{\raggedright\arraybackslash}X|>{\raggedright\arraybackslash}X|c|}
\hline
\rowcolor{TableHeader}
\sffamily\bfseries Tiêu chí & \sffamily\bfseries Xuất sắc (4đ) & \sffamily\bfseries Khá (3đ) & \sffamily\bfseries Cần cải thiện (1-2đ) & \sffamily\bfseries Điểm \\
\hline
\textbf{1. Hiểu mô hình SIR} & Giải thích đúng S,I,R. Tính đúng $R_0$, $H$ & Hiểu cơ bản S,I,R. Tính đúng công thức & Nhầm lẫn khái niệm. Tính sai & \underline{\hspace{1cm}} \\
\hline
\textbf{2. Sử dụng EpidemicSim} & Thành thạo 3 tabs. So sánh 3+ kịch bản & Sử dụng tốt Tab 1,2. So sánh 2 kịch bản & Chỉ dùng Tab 1. Không so sánh được & \underline{\hspace{1cm}} \\
\hline
\textbf{3. Phân tích} & Đề xuất chính sách hợp lý, dựa trên dữ liệu & Đề xuất cơ bản, có lý giải & Không đề xuất được & \underline{\hspace{1cm}} \\
\hline
\textbf{4. Trình bày} & Mạch lạc, có biểu đồ minh họa & Rõ ràng, có hỗ trợ trực quan & Rời rạc, không minh họa & \underline{\hspace{1cm}} \\
\hline
\end{tabularx}

\vspace{0.5cm}
\noindent\textbf{TỔNG ĐIỂM:} \underline{\hspace{2cm}} / 16 điểm

\vspace{0.5cm}
\noindent\textbf{Nhận xét:}

\underline{\hspace{\textwidth}}

\underline{\hspace{\textwidth}}

\vspace{1cm}
%==============================================================
% PHIẾU ĐÁNH GIÁ ĐỒNG CẤP
%==============================================================
\begin{center}
{\Large\bfseries\sffamily\color{primary} PHIẾU ĐÁNH GIÁ ĐỒNG CẤP}
\end{center}

\vspace{0.3cm}
\noindent\textbf{Tên em:} \underline{\hspace{4cm}} \hfill \textbf{Nhóm:} \underline{\hspace{2cm}}

\vspace{0.5cm}
\textbf{Đánh giá các thành viên trong nhóm} (thang 1-5: 1=Yếu, 5=Xuất sắc)

\vspace{0.3cm}
\begin{tabularx}{\textwidth}{|L|c|c|c|c|}
\hline
\header{Tên thành viên} & \header{Ý tưởng} & \header{Hợp tác} & \header{Hoàn thành} & \header{Tổng} \\
\hline
\underline{\hspace{4cm}} & \underline{\hspace{1cm}} & \underline{\hspace{1cm}} & \underline{\hspace{1cm}} & \underline{\hspace{1cm}} \\
\hline
\underline{\hspace{4cm}} & \underline{\hspace{1cm}} & \underline{\hspace{1cm}} & \underline{\hspace{1cm}} & \underline{\hspace{1cm}} \\
\hline
\underline{\hspace{4cm}} & \underline{\hspace{1cm}} & \underline{\hspace{1cm}} & \underline{\hspace{1cm}} & \underline{\hspace{1cm}} \\
\hline
\underline{\hspace{4cm}} & \underline{\hspace{1cm}} & \underline{\hspace{1cm}} & \underline{\hspace{1cm}} & \underline{\hspace{1cm}} \\
\hline
\end{tabularx}

\vspace{0.5cm}
\noindent\textbf{Thành viên đóng góp nhiều nhất:} \underline{\hspace{6cm}}

\end{document}
